\section*{Core equations and notation}

We work in geometrized units \(c=G=1\). For the energy density \(\varepsilon(r)\), pressure \(P(r)\), and enclosed mass \(m(r)\) the Tolman--Oppenheimer--Volkoff equations read
\[
\frac{dm}{dr}=4\pi r^{2}\,\varepsilon(r),
\qquad
\frac{dP}{dr}=-(\varepsilon+P)\,\frac{m+4\pi r^{3}P}{r(r-2m)}.
\]

\subsection*{Variant A: inertial (enthalpy) dressing}
The vacuum-induced response modifies only the inertial (enthalpy) factor:
\[
(\varepsilon+P)\;\longrightarrow\;(\varepsilon+P)\,\bigl[1+I(r;\sigma,x)\bigr],
\]
where \(I(r;\sigma,x)\) is a dimensionless dressing function and \(\sigma\) denotes the scan parameter (we set \(x=1\) throughout the main analysis).

\subsection*{Perturbativity diagnostic}
Define the local perturbativity metric
\[
\mathrm{epsratio}(r)=\frac{\varepsilon_{\mathrm{vac,inertial}}(r)}{\varepsilon_{\mathrm{ref}}(r)}.
\]
Headline results are restricted to
\[
\max_{r}|\mathrm{epsratio}(r)| \le 0.10,
\]
with \(0.10<\max|\mathrm{epsratio}|\le 0.30\) reported as stress tests and \(\max|\mathrm{epsratio}|>0.30\) excluded.

\subsection*{Reactive kernel and effective inertial coefficient}
The low-frequency expansion of the reactive kernel is
\[
K_{R}(\omega)=K_{R}(0)+\tfrac{1}{2}m_{\rm vac}\,\omega^{2}+O(\omega^{4}),
\qquad
m_{\rm vac}=\left.\frac{d^{2}K_{R}}{d\omega^{2}}\right|_{\omega=0}.
\]
We parameterize
\[
m_{\rm vac}=C_{\rm conv}\,\sigma\,x,
\]
and adopt \(C_{\rm conv}=1\), \(x=1\) so that \(\sigma\) is the single scan dial.

\subsection*{Uncertainty propagation}
For an observable \(O(\sigma)\) the matching contribution in the linear regime is
\[
\sigma_{O,\mathrm{match}}=\left|\frac{dO}{d\sigma}\right|\,\sigma_{\mathrm{match}},
\]
and the total uncertainty is combined as
\[
\sigma_{\mathrm{total}}=\sqrt{\sigma_{\mathrm{disc}}^{2}+\sigma_{\mathrm{match}}^{2}+\sigma_{\mathrm{MC}}^{2}}.
\]
If nonlinearity is flagged, linear propagation is replaced by Monte-Carlo envelope propagation through the fitted \(O(\sigma)\) curve and the median with 5--95\% interval is reported.
